% !TeX root = ../../main.tex
\subsection{比較級を使った構文}\vspace{-0.35em}
    \begin{theorembox}{\textbf{比較級の基本構文}}
        \theoremline{A + 比較級 +\fitblankbf[]{than} + B}{「\fitblankbf[]{B} よりも\fitblankbf[]{A}」}
    \end{theorembox}
    \vspace{0.25em}

    ※比較級を使った文を作るときは、基本的に\fitblankbf[]{than}と共に使うことに注意する。

    \begin{techniquebox}{\textbf{比較級を強める語}}
        比較級を強める語として\fitblankbf[]{much}、\fitblankbf[]{by far}、\fitblankbf[]{far}がある。

        これらはいずれも「\fitblankbf[]{ずっと}、\fitblankbf[]{はるかに}」と訳し、比較級の\fitblankbf[]{直前}に置く。
    \end{techniquebox}
    \vspace{0.25em}

    \begin{theorembox}{\textbf{the を使う比較級の構文}}
        \theoremline{\fitblankbf[]{the} + 比較級 +\fitblankbf[]{of the two}}{「\fitblankbf[]{2つのうち}より～な方」}
    \end{theorembox}
    \vspace{0.25em}

    \begin{theorembox}{\textbf{疑問詞を含む比較級の構文}}
        \fitblankbf[]{Which} [\fitblankbf[]{What}] + 動詞 + 比較級\fitblankcomma\fitblankbf[]{A or B} ?
        \theoremmeaning{「\fitblankbf[]{AとB}のうち、\fitblankbf[]{どちらが} (形容詞 [ 副詞 ]) か」}
    \end{theorembox}

\subsection{練習問題}\vspace{-0.35em}
    \begin{enumerate}[leftmargin=*, nosep, itemsep=1.27em]
        \item 彼は私より背が高い。
        \dialogueblankpractice{He is taller than me.}
        \item This problem is easier than that one.
        \dialogueblankpractice{この問題はあの問題より簡単だ。}
        \item 私のかばんはあなたのよりずっと重い。
        \dialogueblankpractice{My bag is much heavier than yours.}
        \item Riku is much smarter than any other person.
        \dialogueblankpractice{リクは他の人よりはるかに頭がいい。}
        \item 2人のうち、彼女のほうが若い。
        \dialogueblankpractice{She is the younger of the two.}
        \item サッカーと野球のうち、どちらがより人気がありますか。
        \dialogueblankpractice{Which is more popular, soccer or baseball?}
    \end{enumerate}

\newpage
